\subsection{Simulación Lineal (Máximo de Prefijo)}

\graybold{Preguntas guía:}

\begin{itemize}
\item ¿La respuesta depende de un solo recorrido lineal del arreglo, calculando un máximo/mínimo/acumulado por posición?
\item ¿No hay dependencias entre posiciones más allá de una comparación directa fila por fila?
\item ¿El problema, aunque tenga una narrativa elaborada, se reduce matemáticamente a una fórmula simple tras identificar la variable clave?
\end{itemize}

\graybold{Utilidad:} Resolver directamente problemas donde, tras modelar bien la situación (aquí: el "cuello de botella" es la fila con menor espacio libre entre dos paredes), un solo recorrido O(n) calculando un máximo/mínimo basta. Es clave en competencia reconocer cuándo un problema de apariencia compleja se reduce a esto.

\graybold{Complejidad:} O(h) tiempo, O(h) espacio (o O(1) extra si se procesa en streaming).

\graybold{Fundamento teórico:} Las dos paredes se mueven a 1 m/s cada una. En la fila i, el espacio libre es w − (l[i] + r[i]) y se cierra a razón de 2 m/s (ambas paredes se acercan). El tiempo en que la fila i se cierra es (w − l[i] − r[i]) / 2. Las paredes chocan en la fila que se cierra primero: la de menor espacio libre inicial, equivalente al máximo de (l[i] + r[i]).

\graybold{Ejercicio aplicado}

Una sala de H filas y ancho W tiene una pared este y una pared oeste que se acercan a 1 m/s cada una desde que alguien entra. Cada fila i tiene un ancho ya ocupado por cada pared. Determinar en cuánto tiempo chocan las paredes.

\graybold{I/O:} H ≤ 10\textasciicircum{}5, dos arreglos de tamaño H → lectura total + tokenización (patrón 2).

\graybold{Código solución}

\begin{lstlisting}[language=Python]
import sys
 
def solve():
    data = sys.stdin.buffer.read().split()
    h = int(data[0]); w = int(data[1])
    left = list(map(int, data[2:2 + h]))
    right = list(map(int, data[2 + h:2 + 2 * h]))
    max_sum = max(left[i] + right[i] for i in range(h))
    ans = (w - max_sum) / 2
\end{lstlisting}


\begin{lstlisting}[language=Python]
if ans.is_integer():        # verifica parseo de float a int
    print(int(ans))
else:
    print(ans)
solve()
\end{lstlisting}